%------------------------------------------------------------------------------%
\documentclass[fleqn,utf8,aspectratio=169,12pt]{beamer}

\usepackage{integracao_e_series}

\title{Séries Numéricas -- Convergência Absoluta}

%------------------------------------------------------------------------------%
\begin{document}

\addtitle

%------------------------------------------------------------------------------%
\section{Convergência Absoluta}

%------------------------------------------------------------------------------%
\begin{frame}{Convergência Absoluta}
  Uma série
  \(\quad \sum_{n=1}^{\infty}\; a_n\) \\[5mm]
  \pause
  \emph{converge absolutamente}, ou é \emph{absolutamente convergente}, \\[5mm]
  \pause
  se
  \(\quad\sum_{n=1}^{\infty} \;\abs{a_n} \quad\) converge
\end{frame}

%------------------------------------------------------------------------------%
\begin{frame}{Convergência Condicional}
  Se  \tabto{12mm} \(\sum_{n=1}^{\infty} a_n\)          \tabto{32mm} \emph{converge} \\[9mm]
  \pause
  mas \tabto{12mm} \(\sum_{n=1}^{\infty} \;\abs{a_n} \) \tabto{32mm} \emph{diverge}  \\[9mm]
  \pause
  a série é \emph{condicionalmente convergente} \\[9mm]
  \pause
  Exemplo: Série Harmônica Alternada
\end{frame}

%------------------------------------------------------------------------------%
\begin{frame}{Teste da Convergência Absoluta}
  Uma série absolutamente convergente é convergente

  \pause
  \vspace{\baselineskip}
  Séries com termos não negativos e convergentes
  são absolutamente convergentes
\end{frame}

%------------------------------------------------------------------------------%
\begin{frame}{Teorema do Rearranjo}
Se \(\quad\sum_{n=1}^{\infty} a_n \quad \) converge absolutamente \\[5mm]

\pause
dado qualquer \emph{rearranjo} $(b_k)$ da sequência $(a_n)$ \\[5mm]

\pause
a série $\sum b_k$ converge e
\[
  \sum_{k=1}^{\infty} b_k \;=\; \sum_{n=1}^{\infty} a_n
\]
\end{frame}

%------------------------------------------------------------------------------%
%\section{Exemplos}

%------------------------------------------------------------------------------%
\section{Lista Mínima}

%------------------------------------------------------------------------------%
\begin{frame}{Lista Mínima}
  Estudar a Seção ? da Apostila

  \vspace{\baselineskip}
  Exercícios:

  \vfill
  Atenção: A prova é baseada no livro, não nas apresentações
\end{frame}

%------------------------------------------------------------------------------%
\end{document}
%------------------------------------------------------------------------------%
